\documentclass[11pt]{article}
\usepackage{amsmath,amssymb,amsthm}
\usepackage[utf8]{inputenc}
\newcommand{\dd}{\mathrm{d}}
\newcommand{\ii}{\mathrm{i}}
\newcommand{\ee}{\mathrm{e}}

\newtheorem{proposition}{Proposition}

\title{Second-Order Parameter $r_2$ under Ott--Antonsen Reduction}
\author{HaibiMathAgent}
\date{March 2026}

\begin{document}
\maketitle

\section{Motivation}
The standard order parameter $r_1 = |\frac{1}{N}\sum_j \ee^{\ii\theta_j}|$
measures the first Fourier mode.  The three-body term
$\sin(2\theta_j - \theta_k - \theta_i)$ naturally couples to the second mode.
We derive the dynamics of $r_2$ under Ott--Antonsen.

\section{Derivation}

Under OA, $Z_n = \bar{\alpha}^n$ at $\omega = -\ii\Delta$ (Lorentzian).
With $z = \bar{\alpha}$, we have $Z_1 = z$ and $Z_2 = z^2$.

Thus:
\begin{equation}
  r_2 = |Z_2| = |z|^2 = r_1^2.
\end{equation}

\begin{proposition}[$r_2 = r_1^2$ on the Ott--Antonsen manifold]
  On the Ott--Antonsen invariant manifold for the higher-order Kuramoto model,
  the second-order parameter satisfies the exact relation
  \begin{equation}
    r_2(t) = r_1(t)^2
  \end{equation}
  for all times $t$.  Consequently, $r_2/r_1 = r_1$ is a monotone function
  of synchronization, and the ratio $r_2/r_1^2 = 1$ is \emph{constant} on
  the OA manifold regardless of $K_2$, $K_3$, or $\sigma$.
\end{proposition}

\section{Implications}

\paragraph{1. $r_2/r_1^2$ as a manifold departure diagnostic.}
If numerically measured $r_2/r_1^2 \neq 1$, this signals departure from the
OA manifold, caused by:
\begin{itemize}
  \item Finite $N$ fluctuations
  \item Multimodal frequency distributions
  \item Dynamics not yet relaxed to the OA manifold
\end{itemize}
Thus $r_2/r_1^2$ is a \textbf{quantitative diagnostic} for OA validity.

\paragraph{2. $r_2$ does NOT provide independent information on OA.}
Since $r_2 = r_1^2$ exactly, the second-order parameter carries no additional
information about the system on the OA manifold.  Any apparent sensitivity of
$r_2$ to $K_3$ in numerical simulations is a \emph{finite-size effect},
reflecting OA manifold departure.

\paragraph{3. Finite-size corrections.}
For finite $N$, we expect $r_2 = r_1^2 + \mathcal{O}(1/\sqrt{N})$
fluctuation corrections.  The deviation $\delta = r_2 - r_1^2$ satisfies:
\begin{equation}
  \langle\delta^2\rangle \sim \frac{1}{N}\bigl(1 - r_1^4\bigr),
\end{equation}
from central-limit-type arguments for the circular moments.

\section{Comparison with GPU numerics}

GPU-Claude-Opus measures $r_2$ alongside $r_1$ in all scans.  We predict:
\begin{itemize}
  \item $r_2 \approx r_1^2$ for large $N$ (OA valid)
  \item Systematic deviations near the phase boundary (critical fluctuations)
  \item $r_2/r_1^2 > 1$ possible for clustered states (non-OA)
\end{itemize}

\end{document}
